% SEGMENT OLP-0564-S01
\documentclass[../../../include/open-logic-section]{subfiles}

\begin{document}

\olfileid{sth}{spine}{rank}
\olsection{தரநிலை}

படிநிலைகளை $V_\alpha$-களாக வரையறுத்துள்ளோம்; மேலும் ஒவ்வொரு கணமும்
ஏதோ ஒரு படிநிலையின் உட்கணம் என்று அறிவோம். ஆகவே ஒரு கணத்தின்
\emph{தரநிலையை} வரையறுக்கலாம். உள்ளுணர்வின்படி, $A$ இன் தரநிலை என்பது
$A$ முதன்முதலில் உருவான தருணம். இன்னும் துல்லியமாக:

\begin{defn}\ollabel{defnsetrank}
ஒவ்வொரு கணம் $A$ க்கும், $A \subseteq V_\alpha$ ஆக அமையும் மீச்சிறு
வரிசையெண் $\alpha$ தான் $\setrank{A}$.
\end{defn}
\begin{prop}\ollabel{ranksexist}
	எந்த $A$ க்கும் $\setrank{A}$ உள்ளது.
\end{prop}
\begin{proof}
	பயிற்சியாக விடப்பட்டுள்ளது.
\end{proof}
\begin{prob}
	\olref[sth][spine][rank]{ranksexist} ஐ நிறுவுக.
\end{prob}
\noindent
தரநிலைகளின் நன்கு வரிசைப்படுத்தல் சில முக்கிய முடிவுகளை நிறுவ
அனுமதிக்கிறது:
\begin{prop}\ollabel{valphalowerrank}
எந்த வரிசையெண் $\alpha$ க்கும்,
$V_\alpha = \Setabs{x}{\setrank{x} \in \alpha}$.
\end{prop}

\begin{proof}
$\setrank{x} \in \alpha$ எனில் $x \subseteq V_{\setrank{x}} \in
V_\alpha$; ஆகவே $V_\alpha$ ஆற்றலுடையது என்பதால்
(\olref[Valphabasic]{Valphabasicprops} ஐப் பலமுறை பயன்படுத்தி)
$x \in V_\alpha$. மறுதிசையில், $x \in V_\alpha$ எனில்
$x \subseteq V_\alpha$; ஆகவே $\setrank{x} \leq \alpha$.
இப்போது ஓர் எளிய முடிவிலிக்கடந்த தொகுத்தறிதல் உண்மையில்
$\setrank{x} \in \alpha$ என்று காட்டுகிறது.
% SOURCE CORRECTION TA-STH-016: replace the contradictory “x is not in V_alpha” conclusion with the required strict rank inequality.
\end{proof}
\begin{prob}
	\olref[sth][spine][rank]{valphalowerrank} இல் குறிப்பிடப்பட்ட எளிய
	முடிவிலிக்கடந்த தொகுத்தறிதலை நிறைவுசெய்க.
\end{prob}

\begin{prop}\ollabel{rankmemberslower}
$B \in A$ எனில், $\setrank{B} \in \setrank{A}$.
\end{prop}

\begin{proof}
\olref{valphalowerrank} இன்படி
$A \subseteq V_{\setrank{A}} = \Setabs{x}{\setrank{x} \in
\setrank{A}}$.
\end{proof}
\noindent
இவ்வுண்மையைப் பயன்படுத்தி, \emph{எல்லாக் கணங்களையும்} பற்றி ஒரு வகைத்
தொகுத்தறிதலால் நிறுவ அனுமதிக்கும் முடிவை அடையலாம்:

\begin{thm}[$\in$-தொகுத்தறிதல் திட்டம்]
எந்த வாய்பாடு $\phi$ க்கும்:
\[
	\forall A((\forall x \in A)\phi(x) \lif \phi(A)) \lif \forall A \phi(A).
\]
\end{thm}

\begin{proof}
எதிர்மறையை நிறுவுவோம். எனவே $\lnot \forall A \phi(A)$ என்க.
முடிவிலிக்கடந்த தொகுத்தறிதலால்
(\olref[ordinals][basic]{ordinductionschema}), சாத்தியமானவற்றுள்
மீச்சிறு தரநிலையுடைய ஏதோ ஓர் $\phi$-அல்லாதது உள்ளது; அதாவது,
$\lnot \phi(A)$ மற்றும் $\forall x(\setrank{x} \in \setrank{A}
\lif \phi(x))$ ஆக அமையும் ஏதோ ஓர் $A$ உள்ளது. இப்போது $x \in A$
எனில், \olref{rankmemberslower} இன்படி $\setrank{x} \in
\setrank{A}$; ஆகவே $\phi(x)$; அதாவது $(\forall x \in A)\phi(x)
\land \lnot \phi(A)$.
\end{proof}\noindent
இவ்வாற்றல்வாய்ந்த முடிவை முறைசாராமல் விளக்கும் ஒரு வழி இதோ. ஒரு கணத்தின்
ஒவ்வொரு !!{element} உம் $\phi$ ஆக இருக்கும்போதெல்லாம் அந்தக் கணமும்
$\phi$ ஆக இருந்தால், $\phi$ \emph{மரபுவழியானது} என்று கூறுக. அப்போது
$\in$-தொகுத்தறிதல் உங்களுக்குப் பின்வருவதைச் சொல்கிறது: $\phi$
மரபுவழியானது எனில், ஒவ்வொரு கணமும் $\phi$ ஆகும்.

தரநிலைகள் பற்றிய விவாதத்தை (இப்போதைக்கு) நிறைவுசெய்ய, சில முறை
முன்னறிவித்துள்ள சில கூற்றுகளை நிறுவுவோம்.

\begin{prop}\ollabel{ranksupstrict}
$\setrank{A} = \supstrict_{x \in A}\setrank{x}$.
\end{prop}

\begin{proof}
$\alpha = \supstrict_{x \in A}\setrank{x}$ என்க.
\olref{rankmemberslower} இன்படி $\alpha \leq \setrank{A}$. ஆனால்
$x \in A$ எனில் $\setrank{x} \in \alpha$; ஆகவே
\olref{valphalowerrank} இன்படி $x \in V_\alpha$, எனவே
$A \subseteq V_\alpha$, அதாவது $\setrank{A} \leq \alpha$.
எனவே $\setrank{A} = \alpha$.
\end{proof}

\begin{cor}\ollabel{ordsetrankalpha}
எந்த வரிசையெண் $\alpha$ க்கும் $\setrank{\alpha} = \alpha$.
\end{cor}

\begin{proof}
எல்லா $\beta \in \alpha$ க்கும் $\setrank{\beta} = \beta$ என்று
முடிவிலிக்கடந்த தொகுத்தறிதலுக்காகக் கொள்க. இப்போது
\olref{ranksupstrict} இன்படி $\setrank{\alpha} =
\supstrict_{\beta \in \alpha}\setrank{\beta} =
\supstrict_{\beta \in \alpha}\beta = \alpha$.
% 	First note that $\setrank{\alpha} \neq \beta$ for any $\beta \in
% 	\alpha$. For otherwise we would have $\beta = \setrank{\beta} \in
% 	\setrank{\alpha} = \beta$ by \olref{rankmemberslower}, i.e.,
% 	$\beta \in \beta$, a contradiction. 
%
% 	Now note that $\alpha \subseteq V_\alpha$. For $\alpha =
% 	\Setabs{\beta}{\beta \in \alpha}$ by
% 	\olref[sfr][ordinals][basic]{ordissetofsmallerord}. And if
% 	$\beta \in \alpha$ then $\beta \subseteq V_{\setrank{\beta}} =
% 	V_\beta \in V_\alpha$, hence $\beta \in V_\alpha$, by
% 	\olref[sfr][spine][Valphabasic]{Valphabasicprops} twice. 
\end{proof}

கடைசியாக, $\ZFminus$ ஆனது
$\emph{ஒழுங்குமை} \Rightarrow \emph{அடித்தளம்}$ என்ற நிபந்தனைவாக்கியத்தை
நிறுவுகிறது என்று \olref[foundation]{sec} இன் இறுதியில் வாக்களிக்கப்பட்ட
முடிவுக்கான விரைவான நிறுவல் இதோ. (``தரநிலை'' என்ற கருத்தும்
\olref{rankmemberslower} உம் இந்த நிறுவலில் பயன்படுத்தக் கிடைக்கின்றன;
ஏனெனில்---இப்பிரிவின் தொடக்கத்தில் குறிப்பிட்டதுபோல---அவற்றை
$\ZFminus + \text{ஒழுங்குமை}$ ஐப் பயன்படுத்தி முன்வைக்கலாம்.)

\begin{prop}[$\ZFminus + \text{ஒழுங்குமை}$ இல் பணிபுரிகிறோம்]\ollabel{zfminusregularityfoundation}
அடித்தளம் நிறைவேறுகிறது.
\end{prop}

\begin{proof}
$A \neq \emptyset$ ஐ நிலைப்படுத்தி, சாத்தியமானவற்றுள் மீச்சிறு தரநிலையுடைய
ஏதோ ஒரு $B \in A$ ஐ எடுத்துக்கொள்க. $c \in B$ எனில்,
\olref{rankmemberslower} இன்படி $\setrank{c} \in \setrank{B}$;
ஆகவே $B$-இன் தேர்வின்படி $c \notin A$.
\end{proof}

\end{document}
